% Generated from locale/tr/content/proof-theory/natural-deduction/translation-G2i.tex; do not hand-edit.


\olfileid[tr]{pt}{nat}{tgi}

\olsection{\Log{G2i}'den \Log{N2i}'ye Çeviri}

\Log{G2i}'deki sekantların öncülleri !!{formula}s{}in çoklu kümeleri~$\Gamma$,
\Log{N2i}'deki sekantların öncülleriyse her etiketin yalnızca bir kez geçtiği
etiketli formül kümeleri~$\Gamma'$'dür. Etiketli !!{formula}s{}in bir
$\Gamma'$ kümesinin, her $!A$ !!{formula}{}ü için $\Gamma'$ içinde $x : !A$
biçimindeki !!{element} sayısı $\Gamma$ içindeki $!A$'nın katlılığından (yani
$!A$ kopyalarının sayısından) küçük ya da ona eşitse $\Gamma$ çoklu kümesine
\emph{karşılık geldiğini} söyleyelim.

\begin{prop}\ollabel{prop:G2i-to-N2i}
$\Log{G2i} + \Cut \Proves \Gamma \Sequent \Delta$ ise $\Log{N2i}
\Proves \Gamma' \Sequent !C$ olur; burada $\Delta$ boşsa $!C \ident \lfalse$,
aksi hâlde $\Delta = \{!A\}$'dır ve $\Gamma'$, $\Gamma$'ya karşılık gelir.
\end{prop}

\begin{proof}
  $\pi$, $\Log{G2i}+\Cut$ içinde $\Gamma \Sequent \Delta$'nın bir !!a{proof}{}ı
  ise $\Gamma$'ya karşılık gelen bir~$\Gamma'$ ile \Log{N2i} içinde
  $\Gamma' \Sequent !C$'nin bir~$\delta$ !!a{proof}{}ı bulunduğunu gösteriyoruz.
  Bunu $n = \pheight{\pi}$ üzerinde tümevarımla yapıyoruz.

  $n = 0$ ise $\pi$ bir aksiyomdur. $\lfalse \Sequent \quad$ aksiyomuysa
  $\delta$, $x:\lfalse \Sequent \lfalse$ aksiyomudur; yani
  $\Gamma' = \{x:\lfalse\}$ ve $!C \ident \lfalse$ olur. $\pi$, $!A \Sequent
  !A$ biçiminde bir aksiyomsa $\delta$, $x: !A \Sequent !A$'dır.
  
  $n>0$ ise $\pi$ bir $R$ kuralıyla biter. Tümevarım varsayımı, bu kuralın
  öncüllerine karşılık gelen sekantların \Log{N2i} içindeki !!{proof}s{}ını
  güvence altına alır. Bu ispatlardaki tüm boşaltma etiketlerinin birbirinden
  farklı olduğunu ve bir $x$ etiketinin, boşaltılıyorsa yalnızca onu boşaltan
  çıkarımın üstünde geçtiğini (gerekirse yeniden etiketleyerek) varsayabiliriz.
  Ardından bu !!{proof}s{}ın uygun bir $\Gamma' \Sequent !C$ ispatı sağlayacak
  biçimde nasıl birleştirilebileceğini gösteririz. $R$'ye göre durumları
  ayırıyoruz.

  \begin{enumerate}
    \item $R$, \RightR{\Weakening} ise $\pi$ şöyle biter:
    \[
    \AxiomC{}
    \RightLabel{$\pi_1$}
    \Deduce$\Gamma \fCenter$
    \RightLabel{\RightR{\Weakening}}
    \UnaryInf$\Gamma \fCenter !A$
    \DisplayProof
    \]
    $\pi_1$'e uygulanan tümevarım varsayımı $\Gamma' \Sequent \lfalse$'ın bir
    !!a{proof}{}ını verir. $\Gamma' \Sequent !A$'nın bir !!a{proof}{}ını elde
    etmek için $\FalseInt$ uygularız.
    \item $R$, \LeftR{\Weakening} ise $\pi$ şöyle biter:
    \[
    \AxiomC{}
    \RightLabel{$\pi_1$}
    \Deduce$\Gamma \fCenter \Delta$
    \RightLabel{\LeftR{\Weakening}}
    \UnaryInf$!A, \Gamma \fCenter \Delta$
    \DisplayProof
    \]
    $\pi_1$'e uygulanan tümevarım varsayımı $\Gamma' \Sequent !C$'nin bir
    !!a{proof}{}ını verir. Bunu $\delta$ olarak alırız. $\Gamma'$, $\Gamma$'ya
    karşılık geliyorsa $!A, \Gamma$'ya da karşılık geldiğine dikkat edin:
    $\Gamma'$ içindeki $x : !A$ sayısı $\Gamma$ içindeki $!A$ kopyalarının
    sayısından büyük değildir; dolayısıyla $\Gamma \cup \{!A\}$ içindeki
    $!A$ kopyalarının sayısından da büyük değildir.
    \item $R$, \LeftR{\Contraction} ise $\pi$ şöyle biter:
    \[
    \AxiomC{}
    \RightLabel{$\pi_1$}
    \Deduce$!A, !A, \Gamma \fCenter \Delta$
    \RightLabel{\LeftR{\Contraction}}
    \UnaryInf$!A, \Gamma \fCenter \Delta$
    \DisplayProof
    \]
    $\pi_1$'e uygulanan tümevarım varsayımı $\Pi, \Gamma' \Sequent !C$'nin
    bir $\delta_1$ !!a{proof}{}ını verir; burada
    $\Pi \subseteq \{x : !A, y:!A\}$'dır.
    $\Pi = \{x : !A, y:!A\}$ ise $\delta_1$ içindeki bütün etiketli
    $y:!A$ !!{formula}{}ünü $x:!A$ ile değiştirin. $x$, $\delta_1$'in son
    sekantının bağlamında geçtiğinden $\delta_1$ içinde hiçbir çıkarımın
    $x:!B$ biçimindeki bir !!a{formula}{}ü boşaltmadığını, yani $\delta_1$
    içindeki hiçbir sekantın solundaki $x:!B$'nin etkin olmadığını
    varsayabiliriz. Dolayısıyla sonuç hâlâ \Log{N2i} içinde bir !!a{proof}{}tır.
    \item $R$, $\RightR{\lif}$ ise $\pi$ şöyle biter:
    \[
    \AxiomC{}
    \RightLabel{$\pi_1$}
    \Deduce$!A, \Gamma \fCenter !B$
    \RightLabel{\RightR{\lif}}
    \UnaryInf$\Gamma \fCenter !A \lif !B$
    \DisplayProof
    \]
    Tümevarım varsayımı gereği $x:!A, \Gamma' \Sequent !B$'nin ya da
    $\Gamma' \Sequent !B$'nin bir $\delta_1$ !!a{proof}{}ı vardır; burada
    $\Gamma'$, $\Gamma$'ya karşılık gelir. $\delta$ olarak $\delta_1$'in
    sonuna bir \Intro{\lif} uygulaması eklenmiş biçimini alabiliriz.
    \item $R$, \LeftR{\lif} ise $\pi$ şöyle biter:
    \[
    \AxiomC{}
    \RightLabel{$\pi_1$}
    \Deduce$\Gamma_1 \fCenter !A$
    \AxiomC{}
    \RightLabel{$\pi_2$}
    \Deduce$!B, \Gamma_2 \fCenter \Delta$
    \RightLabel{\LeftR{\lif}}
    \BinaryInf$!A \lif !B, \Gamma_1, \Gamma_2 \fCenter \Delta$
    \DisplayProof
    \]
    Tümevarım varsayımı !!{proof}s $\delta_1$ ile $\delta_2$'yi verir:
    birincisi $\Gamma_1' \Sequent !A$'yı, ikincisi ya $x: !B, \Gamma_2'
    \Sequent !C$'yi ya da $\Gamma_2' \Sequent !C$'yi ispatlar.
    $\delta_2$ ikinci biçimdeyse $\delta_2$, $\delta$ olarak kullanılabilir.
    Aksi hâlde hiçbir etiketin hem $\delta_1$ hem de $\delta_2$ içinde
    geçmemesini sağlamak için $\delta_1$ içindeki bütün formülleri ve boşaltma
    etiketlerini yeniden etiketleyebiliriz. Bu durumda $\Gamma_1' \cup
    \Gamma_2'$, $\Gamma_1 \cup \Gamma_2$'ye karşılık gelir. $x:!B$,
    $\delta_2$'nin son sekantında geçtiğinden $\delta_2$ içinde $x$ boşaltma
    etiketli hiçbir çıkarımda etkin olamaz. $\delta_2$ içindeki her
    $x:!B \Sequent !B$ aksiyomunu şu $\delta_3$ !!{proof}{}ıyla değiştirin:
    \[
    \Axiom$x: !A \lif !B \fCenter !A \lif !B$
    \AxiomC{}
    \RightLabel{$\delta_1$}
    \Deduce$\Gamma_1' \fCenter !A$
    \RightLabel{\Elim{\lif}}
    \BinaryInf$x: !A \lif !B \fCenter !B$
    \DisplayProof
    \]
    ve $\delta_2$ içindeki her $x:!B$ geçişini $x:!A \lif !B$ ile değiştirin.
    Sonuç $\Subst{\delta_2}{\delta_3}{x:!B}$, $x:!A \lif !B, \Gamma_1',
    \Gamma_2' \Sequent !C$'nin bir !!a{proof}{}ıdır.
    \item $R$, \RightR{\land} ise $\pi$ şöyle biter:
    \[
    \AxiomC{}
    \RightLabel{$\pi_1$}
    \Deduce$\Gamma_1 \fCenter !A$
    \AxiomC{}
    \RightLabel{$\pi_2$}
    \Deduce$\Gamma_2 \fCenter !B$
    \RightLabel{\RightR{\land}}
    \BinaryInf$\Gamma_1, \Gamma_2 \fCenter !A \land !B$
    \DisplayProof
    \]
    Tümevarım varsayımı, $\Gamma_1' \Sequent !A$'nın $\delta_1$ ispatını
    ve $\Gamma_2' \Sequent !B$'nin $\delta_2$ ispatını verir.
    $\delta_1$ ve $\delta_2$ içinde hiçbir etiketin ortak geçmemesini sağlamak
    için $\delta_2$ içindeki tüm formülleri ve boşaltma etiketlerini yeniden
    etiketleyebiliriz. Bu durumda $\Gamma_1' \cup \Gamma_2'$, $\Gamma_1 \cup
    \Gamma_2$'ye karşılık gelir. $\delta_1$ ile $\delta_2$'nin son sekantlarına
    $\Intro{\land}$ uygulayarak $\delta$ !!a{proof}{}ını elde ederiz.
    \item $R$, $\LeftR{\land}$ ise $\pi$, örneğin şöyle biter:
    \[
    \AxiomC{}
    \RightLabel{$\pi_1$}
    \Deduce$!A, \Gamma \fCenter \Delta$
    \RightLabel{\LeftR{\land}}
    \UnaryInf$!A \land !B, \Gamma \fCenter \Delta$
    \DisplayProof
    \]
    Tümevarım varsayımı gereği $x:!A, \Gamma' \Sequent !C$'nin ya da
    $\Gamma' \Sequent !C$'nin bir $\delta_1$ !!a{proof}{}ı vardır; burada
    $\Gamma'$, $\Gamma$'ya karşılık gelir. İkinci durumda $\delta_1$'i
    $\delta$ olarak alabiliriz. Birinci durumda $x:!A$ son sekantta geçtiği
    için $\delta_1$ içindeki $x$ boşaltma etiketli hiçbir çıkarımda etkin
    olmadığına dikkat edin. Her $x:!A \Sequent !A$ aksiyomunu şu $\delta_2$
    !!{proof}{}ıyla değiştirin:
    \[
    \Axiom$x: !A \land !B \fCenter !A \land !B$
    \RightLabel{\Elim{\land}}
    \UnaryInf$x: !A \land !B \fCenter !A$
    \DisplayProof
    \]
    ve her $x:!A$ geçişini $x:!A \land !B$ ile değiştirin; yani
    $\Subst{\delta_1}{\delta_2}{x:!A}$'yı ele alın. Bu, $x:!A \land !B,
    \Gamma' \Sequent !C$'nin bir !!a{proof}{}ıdır.
    \item $R$, \Cut ise $\pi$ şöyle biter:
    \[
    \AxiomC{}
    \RightLabel{$\pi_1$}
    \Deduce$\Gamma_1 \fCenter !A$
    \AxiomC{}
    \RightLabel{$\pi_2$}
    \Deduce$!A, \Gamma_2 \fCenter \Delta$
    \RightLabel{\Cut}
    \BinaryInf$\Gamma_1, \Gamma_2 \fCenter \Delta$
    \DisplayProof
    \]
    Tümevarım varsayımı $\Gamma_1' \Sequent !A$'nın $\delta_1$ ispatını
    ve ya $x:!A, \Gamma_2' \Sequent !C$'nin ya da $\Gamma_2' \Sequent !C$'nin
    $\delta_2$ ispatını verir. İkinci durumda $\delta$ olarak
    $\delta_2$'yi alın. Aksi hâlde $\delta_2$ içindeki her
    $x:!A \Sequent !A$ aksiyomunu $\delta_1$ ile değiştirin; böylece
    $\Gamma_1', \Gamma_2' \Sequent !C$'nin
    $\Subst{\delta_2}{\delta_1}{x:!A}$ !!{proof}{}ı olarak~$\delta$'yı elde
    edersiniz.
  \end{enumerate}
\end{proof}

